% !TeX root = ../fundamentals.tex
% Section 2.4 -- Datasets and evaluation
% Draft 1 (2026-07-21). Obeys WRITING_LAW.md + GLOSSARY.md. Citation/number ledger at the section end.

\section{Datasets and evaluation}
\label{sec:fund:eval}

A claim about mobility prediction is only as trustworthy as the data it is measured on
and the protocol that measures it. This section fixes both: the datasets the
dissertation uses, the metrics and reference points each result is read against, the
validation protocol that keeps the estimate honest, and the tests that license the
verbs used to report a comparison.

Two check-in datasets serve as the ground. Gowalla is the dataset of record,
introduced with the study of periodic and socially driven movement in LBSNs, and this
work uses five United States states from it \cite{cho2011gowalla}. For a
non-United-States setting, the Massive-STEPS benchmark supplies Istanbul check-ins and
argues that the field has leaned too long on decade-old datasets, which is part of the
reason for evaluating across more than one city \cite{wongso2025massivesteps}. The
Foursquare check-in dataset from New York and Tokyo is a further widely used benchmark
\cite{yang2015tsmc}; it is named here for context, as the dissertation evaluates on
Gowalla and the Istanbul data rather than on it.

The category label distribution is imbalanced. The Food class alone accounts for
roughly a third of the check-ins in a representative state, so plain accuracy is
inflated by the majority class and must not be the headline for next category.
Macro-averaged F1 is the primary category metric instead: it is the unweighted mean of
the per-class F1 scores, so it reads as average per-class quality with every class
counting equally, and it does not reward a model that serves only the frequent classes
\cite{sokolova2009measures}. The training pipeline counters the same imbalance with
class-weighted cross-entropy. For next region, the primary metric is accuracy at 10
(Acc@10), the fraction of cases where the true region appears in the model's ten
highest-ranked predictions; it is a ranking metric and says nothing about the
probability mass placed on the true region, which is why mean reciprocal rank
accompanies it where the joint comparison needs a rank-sensitive figure. In the
reported region results a region never seen in training counts as a miss, so the
figure is not inflated by unpredictable new regions; Chapter~5 gives the full
definition. When one model must be scored against its single-task
references across both tasks at once, the aggregate is the relative multi-task
performance change, the average per-task percentage by which the joint model leads or
trails the dedicated single-task models \cite{maninis2019attentive}. It is a relative
figure against a named baseline, and a positive value is a lead only when each
per-task change is itself established.

Every metric is read against reference points. The majority-class floor, which always
predicts the most frequent category, is the level a learned category model must clear.
A first-order transition model, a mobility Markov chain over the training partition,
is the corresponding non-learned floor for the sequential targets \cite{gambs2012mmc}.
The entropy analysis of human movement that reports a potential predictability of about
93\% sets a bound for predicting the next location at coarse resolution, and it frames
why mobility is learnable at all \cite{song2010limits}; it is not, however, a ceiling on
seven-class category macro-F1 or on region ranking, which are different label spaces. The
operative ceiling for the two tasks studied here is the dedicated single-task model,
since the question throughout is whether one joint model can match or outperform the
models built for each task alone.

The validation protocol guards against the most damaging error in this setting, a user
whose check-ins appear in both training and test. Estimates use stratified k-fold
cross-validation \cite{kohavi1995crossval}, and the folds are formed so that no user
spans a split: a grouped, stratified splitter keeps all of a user's check-ins on one
side of every fold, so that measured accuracy reflects generalization to new users
rather than memorization of familiar ones \cite{pedregosa2011sklearn}. Comparisons
are then made across the folds and repetitions with paired tests. The paired Wilcoxon
signed-rank test compares two models across the paired results without assuming
normality, and it is the test that licenses the verb ``outperforms''
\cite{wilcoxon1945}; the Holm step-down procedure controls the family-wise error when
several datasets are tested at once \cite{holm1979}. A superiority test cannot, by
itself, support a claim of equivalence, so where the dissertation states that one
model matches another it uses two one-sided tests to establish statistical
non-inferiority within a two-point margin \cite{lakens2017tost}. The verbs and the
tests are bound together throughout the document: ``outperforms'' follows only from a
paired superiority test, and ``matches'' only from a non-inferiority test within the
stated margin.

% ---------------------------------------------------------------------------
% CITATION / NUMBER LEDGER -- 2.4 (numbers quoted, not computed; conventions named)
%  DATASETS FIX (fact F4): "two ... serve as the ground" = Gowalla + Massive-STEPS/Istanbul (the two USED).
%                        Foursquare (yang2015tsmc) is named as CONTEXT ONLY, explicitly not evaluated on (was mis-counted as a 3rd "ground").
%  cho2011gowalla        Gowalla origin, 5 US states used (KDD 2011)
%  wongso2025massivesteps Massive-STEPS Istanbul benchmark; over-reliance on old datasets [preprint arXiv:2505.11239]
%  yang2015tsmc          Foursquare NYC/Tokyo (TSMC2014), most-cited (815), covers both cities. CONTEXT ONLY (not used). [new bib]
%  sokolova2009measures  macro-F1 = mean of per-class F1 (VERIFIED PDF). Primary category metric. [new bib]
%  maninis2019attentive  Delta_m = avg relative per-task change vs single-task (VERIFIED PDF, Sec.4 + Table 3). Sign convention defined in addendum. [new bib]
%  gambs2012mmc          mobility Markov chain = non-learned floor (targets next PLACE) [new bib]
%  song2010limits        93% = next-LOCATION predictability bound (VERIFIED firsthand PDF). FIX (both gates F5): NOT presented
%                        as a ceiling on category macro-F1 or region Acc@10; the dedicated single-task model is the operative ceiling.
%  kohavi1995crossval    stratified k-fold CV (VERIFIED by author). [new bib; Zenodo DOI 10.5281/zenodo.19712698 for a resolvable id]
%  pedregosa2011sklearn  scikit-learn + StratifiedGroupKFold user-disjoint splitter. SINGLE cite (author ruling); splitter behavior stated defensively (v1.0/2021 addition). [new bib]
%  wilcoxon1945          paired signed-rank test ("outperforms" bound to it) (VERIFIED PDF). [new bib]
%  holm1979              Holm step-down family-wise correction
%  lakens2017tost        TOST equivalence within a margin ("matches" bound to it, two-point margin) [new bib]
%  NUMBERS: 93% (song2010limits, verified, scope-corrected); "roughly a third" Food = a REPRESENTATIVE STATE (Alabama
%           Next-Category CHECK-IN distribution, Food 34.2% per docs/context/TASKS.md; qualitative, NOT recomputed;
%           fact F9: addendum provenance pointer to be corrected to the check-in table, not the 32.5% POI-count table);
%           "10" in Acc@10 (metric definition); "two-point margin" (TOST convention, GLOSSARY).
%  OOD FORESHADOW (domain F8): region metric foreshadowed as "unseen region counts as a miss"; full OOD-Acc@10 definition in Ch.5.
%  Acc@10, MRR, majority-class floor: conventions with no origin DOI -- defined defensively, cited conventionally.
%  pedregosa2011sklearn: StratifiedGroupKFold is a scikit-learn v1.0/2021 feature; single-cite ruling (author) -- 2011 paper for the
%           library, splitter behavior stated defensively in prose. kohavi1995crossval: claim PLAUSIBLE (Zenodo re-deposit id), author to confirm original IJCAI-95 text.
%  Delta_m written in prose as "relative multi-task performance change"; formula lives in the metrics addendum, not inline here.
%  four-of-six + two-point margin (referenced in 2.5): confirm against Ch.5 RESULTS_BOARD.md / PAPER_PLAN §3 at adaptation (fact NOTE 10).
% ---------------------------------------------------------------------------
